\section{Literature Review, Research Gaps, and Theoretical Framework}

\subsection{Gendered Space in Campus Gyms}

Campus gyms are not neutral containers but spatial practices continuously shaped by social relations.
Research on women's entry into gyms has repeatedly shown that the weight area is not merely delineated by equipment---it is simultaneously encoded with gendered meanings:
the concentration of equipment, the high intensity of movements, the prominent noise, the extended equipment occupancy times, and the greater visibility of male users all contribute to its being experienced as a ``male space.''
\textcite{turnock2021gyms}'s qualitative study of commercial gyms makes this point especially clear:
women's barriers to entry arise not only from whether ``women know how to train,'' but from the sharp gendered segregation of the weight area, the emotional threshold of crossing into that space, the performance of masculinity within it, and the discomfort of being continuously exposed to others' gazes.
\textcite{rendall2024gym}'s recent study of female college students further demonstrates that this spatial deterrence is not an isolated narrative but a measurable, significant phenomenon within their campus female sample.
These findings establish that gendered spatial encoding is a precondition for understanding why women's training participation takes specific avoidance forms---a question to which the resistance training literature now turns.

\subsection{Barriers to Resistance Training}

Alongside spatial barriers are participation barriers specific to resistance training itself.
\textcite{peters2019barriers} surveyed college-aged women and identified resistance training barriers as a distinct phenomenon rather than a subcomponent of general physical inactivity: women reported lack of knowledge about proper technique, uncertainty about how to design a training program, perceived social judgment in the weight room, and limited access to guidance or instruction.
These barriers are compounded by the gendered spatial dynamics discussed above: when the free-weight area is already experienced as a ``male space,'' the absence of equipment knowledge and instructional support makes entry even more daunting.
\textcite{rendall2024gym}'s finding that gym intimidation---measured through the SEAM subscale \parencite{levinson2013seam}---functions as a significant exercise barrier among female college students further supports the view that spatial and knowledge-based barriers are intertwined rather than independent.
For this reason, the present study treats ``resistance training avoidance'' rather than ``insufficient physical activity'' as the more precise behavioral outcome variable, asking why a considerable proportion of women concentrate their exercise practices on cardio and light toning while avoiding free weights.
These participation barriers, however, do not arise in a vacuum; they are experienced through gendered bodies already shaped by processes of objectification and surveillance.

\subsection{Objectification and Body Surveillance}

If spatial theory helps us understand ``where entry becomes difficult,'' theories of the female body help us understand ``why the same space is more hostile to certain bodies.''
This study maintains its theoretical anchoring in \textcite{lefebvre1991production} and \textcite{beauvoir1949second}:
the former helps us understand that space is not a neutral given but a product of power relations; the latter frames the female body as an experiential process of being shaped, evaluated, and objectified under the social gaze.
At the level of empirical operationalization, this line of inquiry aligns closely with \textcite{fredrickson1997objectification}'s objectification theory.
This theory posits that women's tendency to treat their own bodies as objects to be scrutinized is not merely an individual psychological deficit but is tied to enduring sociocultural structures of gendered looking.
In gym environments---characterized by abundant mirrors, open sight lines, movements that can easily be performed ``incorrectly,'' and the high visibility of others' bodily displays---this body surveillance becomes more likely to translate into shame, anxiety, and avoidance.
Crucially, the body surveillance triggered in gym spaces does not operate in isolation; it is amplified by an increasingly pervasive digital culture of appearance norms, which the next section examines.

\subsection{Social Media Beauty Norms and Internalization}

Social media appearance norms provide a continuous cultural backdrop for this spatial anxiety.
Relevant reviews and systematic reviews have converged on the finding that social network use is consistently associated with body image disturbance, disordered eating outcomes, and appearance-based self-surveillance, with appearance comparison and appearance ideal internalization serving as key mechanisms \parencite{holland2016systematic}.
\textcite{jung2022social}'s study of young women further indicates that the effect of social media use intensity on body image is primarily mediated by appearance ideal internalization and social comparison, rather than being a simple time-exposure effect.
For young women today, what warrants attention is no longer merely the traditional thin ideal, but the fit ideal that simultaneously emphasizes ``slimness,'' ``firmness,'' ``self-discipline,'' and ``health as display.''
\textcite{donovan2020fitideal} thus argue that ``strong is the new skinny'' does not mean the disappearance of appearance pressure; rather, the pressure has shifted from ``must be thin'' to ``must be strong in the right way.''
This finding is highly consistent with \textcite{tiggemann2015fitspiration}'s experimental evidence:
fitspiration images that appear positively motivated may nonetheless undermine women's body satisfaction in a short period of time.
These findings, however, are predominantly situated within Western media ecosystems; whether and how they manifest in East Asian digital cultures---where beauty norms, fitness discourses, and body ideals may follow different logics---requires further contextualization.

\subsection{The East Asian and Chinese Context}

Placing this research chain within the East Asian and Chinese context renders the issues even more concrete.
Chinese college student research has shown that media and peers influence negative body image through appearance comparison and thin ideal internalization \parencite{shen2022media}, and that Chinese young women's social media practices are linked to body surveillance, weight-control exercise, and restrictive eating:
\textcite{yao2020selfie} found that selfie posting and body surveillance promote a tendency toward ``exercising to control weight'';
another study of Chinese female college students showed that SNS body comparison can influence restrictive eating through body shame \parencite{yao2021bodyimage}.
In other words, in Chinese samples, social media does not merely ``make people unhappy''---it may also shape how people understand exercise, why they exercise, and whether they pursue bodily capability or appearance management.
However, this process should not be understood as one-dimensional passive compliance.
Through the fitness narratives of 12 Chinese women, \textcite{xiong2021fanshen} situates dieting, weight loss, body management, fitness sociality, body stereotypes, self-identity, and cultural aesthetics within a single experiential landscape, suggesting that women's fitness may be organized by appearance norms but may also become part of life planning, self-understanding, and bodily capability reconstruction.
Thus, women may oscillate among discipline, sociality, self-formation, emotional recovery, and agentic negotiation.
This agentic dimension is important because it suggests that women are not simply passive recipients of spatial and appearance pressures---yet whether and how such agency operates specifically within campus gym free-weight areas remains underexplored.

\subsection{Research Gaps}

The existing literature has thus separately addressed gendered space in gyms, social media body image, and appearance pressure among Chinese college students, yet a clear intersectional gap remains:
First, research on women's gym entry barriers has predominantly focused on European and American commercial gyms or general female samples, with limited direct attention to the more specific context of ``Asian female college students + campus gyms + free-weight area.''
Second, research on Chinese college students' body image and social media has primarily focused on negative body image, body shame, restrictive eating, and selfie behavior, with less attention to connecting these psychological and cultural mechanisms to specific training choices within campus fitness spaces.
Third, spatial pressure and digital appearance pressure have typically been studied separately rather than examined simultaneously within a mixed methods framework.
The value of the present study lies precisely in bringing these two usually separate questions together.

\subsection{Integrated Theoretical Framework}

Drawing on the above literature, this study proposes an integrated theoretical framework.
The first layer is \textbf{spatial production logic}: as demonstrated by \textcite{turnock2021gyms} and \textcite{rendall2024gym}, the free-weight area in campus gyms is not a naturally formed training venue but is continuously constructed as a male-dominated area through its layout, equipment, user composition, social norms, and bodily display; this spatial encoding translates into contextualized social exercise anxiety and gym avoidance.
The second layer is \textbf{body objectification and appearance internalization logic}: following \textcite{fredrickson1997objectification} and extending the media effects literature reviewed above \parencite{holland2016systematic,jung2022social}, women are more likely to engage in body surveillance in image-based media and public exercise environments, transforming external appearance standards into self-imposed requirements.
The third layer is the \textbf{behavioral outcome layer}: as suggested by \textcite{peters2019barriers}'s identification of resistance training as a distinct behavioral domain and \textcite{xiong2021fanshen}'s account of agentic negotiation in women's fitness, when ``this space does not belong to me'' and ``my body must not deviate from the ideal femininity'' co-occur, women are more likely to concentrate their exercise choices on cardio, fat loss, light toning, or peripheral machine areas, while exhibiting avoidance, postponement, strategic negotiation, or conditional entry with respect to free-weight training.

This integrated framework also naturally corresponds to the present study's measurement design:
SEAM is suited to capturing contextualized social exercise anxiety and avoidance; SATAQ-4 provides a measurement reference for understanding media pressure and appearance ideal internalization; and semi-structured interviews are used to complement the narrative dimension of ``how they explain, experience, and negotiate these pressures.''
